\section{Introduction}
A computational policy is useful only to the extent that its approximation errors can be related to the decisions and welfare objects of the economic model. This requirement becomes demanding when controls affect both the distribution of future resources and the agent's preferences. An error in the continuation value can change consumption, portfolio risk, and the incentive to adjust preferences. A numerically feasible policy can also solve the wrong economy if the treatment of a state-space boundary implicitly adds resources. For recursive preferences and strategic interaction, even the object called a continuation value must be specified before an algorithm can be evaluated.

This paper develops Neural Bellman Operators (NBO) as a separated approximation of economic control operators. A critic represents a policy's continuation value. An actor proposes a feasible control using that continuation value. A distinct improvement calculation checks the proposal, and an independent evaluation measures the value of the resulting policy. We use both differential Hamiltonian operators and finite-step transition operators. The latter provide an executable route from a specified stochastic economy to neural training without requiring the critic's second derivatives. Their approximation of a continuous-time model is assessed separately from optimization of the finite-step problem.

The distinction between proposal and improvement is important. A stationary point of a neural actor's parameter objective need not maximize the Hamiltonian. An exact pointwise policy-improvement theorem therefore does not establish convergence of an unconstrained neural optimizer. We address this difference constructively. A candidate safeguard includes the actor's proposal in a feasible comparison set. For convex control problems, a second safeguard uses a computable linear minimization gap to bound the remaining improvement error. Neither safeguard assumes that a small parameter gradient implies an economically optimal action.

Our first result gives an error account for finite-horizon policy value. It applies to monotone Bellman operators and permits a different approximation error at each date. The bound distinguishes the critic's evaluation defect, the actor's improvement gap, terminal or boundary data, and the discrepancy between the implemented and target transition operators. The theorem concerns functions and operators, not a particular optimizer. It becomes a numerical certificate only when its error bounds have been established on the required domain. This distinction allows ordinary held-out diagnostics and rigorous finite-problem certificates to appear in the same computational framework without being confused.

Our second result concerns endogenous preferences. In the portfolio economy, risk aversion is a state variable that can be adjusted at a quadratic cost. We specify a liquidation rule at the first exit from the operating domain. The model retains stochastic preference shocks, financial risk, their covariance, and all three controls. It does not replenish wealth when a simulated path crosses a lower boundary. For this economy, the value is the upper envelope of affine functions of the adjustment-cost coefficient. Consequently, the value is decreasing and convex in that coefficient, and optimal discounted cumulative adjustment effort is decreasing. The conclusion incorporates changes in portfolios, consumption, and the stopping time. It does not require, and does not imply, pointwise monotonicity of the preference-adjustment policy.

We implement the separated operator in two economically different environments. The portfolio exercise uses ordinary multilayer neural networks and a finite-candidate improvement safeguard. Its independent evaluation exposes substantial differences between initial-state performance and uniform accuracy near the liquidation boundary. A second application allocates investment across interacting stocks subject to a shared resource budget. It has correlated shocks, a nonseparable state cost, and binding constraints. A convex neural critic supports globally checked one-period improvement, while an independent nonanticipative scenario-tree optimization bounds policy cost losses at the specified initial states. This experiment supplies a dynamic, constrained comparison; it is not a static quadratic exercise relabeled as a scalability test.

The paper also retains recursive utility, temporal selves, dynamic games, nonsmooth solutions, and stochastic trace calculations. These extensions play different roles. Homothetic Merton and Epstein--Zin problems test actual evaluation and actor updates against independently derived policies. A sophisticated consumption problem tests the distinction between the current self's criterion and ordinary continuation evaluation. A stochastic capacity game tests unilateral dynamic deviations, not just joint-profit maximization. Finally, counterexamples and trace calculations identify invalid shortcuts in neural approximation. These applications are not pooled into a single numerical success rate.

\subsection{Related Literature and the Contribution}
Policy improvement precedes neural approximation. \citet{jacka2015} give continuous-time policy-improvement results under explicit stochastic-control hypotheses. Our exact-operator proposition is a conditional statement within that tradition; its purpose is to make the limiting argument and boundary requirements explicit. The new computational analysis concerns the errors that remain when evaluation and improvement are approximate.

The closest neural comparisons also combine value approximation and control improvement. \citet{kim2025} use neural fixed-policy PDE evaluation within policy iteration. Their subsequent revision, \citet{kim2026}, develops interior error bounds and analyzes unresolved boundary information. \citet{lee2024} combine policy iteration with deep operator learning, including transfer across terminal functions. \citet{cai2025} train value and control representations using martingale conditions without requiring an explicit Hamiltonian optimizer. Direct approximation of stochastic controls by neural networks is already present in \citet{han2016}. Thus neural evaluation, policy factorization, mesh-free training, and avoidance of a closed-form argmax are not claims of priority in this paper.

Our contribution is the combination of an explicit economic boundary contract, a value-loss account for the implemented operator, and improvement safeguards whose error has a direct interpretation. The finite-horizon bound is deliberately architecture-independent. The constrained stochastic experiment instantiates its actor component with a convex certificate, and the endogenous-preference result supplies an economic comparative static that does not rely on an unverified policy surface. These are specific claims rather than an assertion that a general combination of neural networks and dynamic programming is new.

The broader numerical literature includes projection and recursive methods \citep{judd1998,stokey1989}, sparse-grid approaches \citep{brumm2017}, and neural differential-equation solvers \citep{raissi2019,han2018}. The relevant comparison depends on structure. A problem with an exact Riccati representation should be compared with that representation. A constrained convex dynamic program should be compared with a solver that exploits convexity. A general nonlinear problem requires its own reference or verification argument. We do not infer a universal dimension threshold from the cost of differentiating a network.

Convex neural parameterizations and optimization over their inputs have an established role in control; see \citet{amos2017}. Our resource-allocation implementation uses a transparent convex feature class rather than proposing a new input-convex architecture. Its role here is to make the feasible improvement gap computable and to connect it to an independently bounded policy cost.


The following sections define the economic operator, establish approximation results, develop the preference economy, and report the computations. The final substantive section treats recursive utility, temporal selves, and dynamic games. A standalone supplement provides proofs, implementation details, additional diagnostics, and the conditions needed to interpret extensions beyond the executed experiments.
